\documentclass{article}
\usepackage{../math_template}

\author{Jonathan Kasongo}
\title{Simon Marais 2025 --- B1}

\begin{document}
\maketitle

\begin{problem}{Simon Marais 2025 --- B1}
A function $f:\{1,2, \ldots, n\} \rightarrow\{m+1, m+2, \ldots, m+n\}$,
where $m$ and $n$ are positive integers, is said to be \textit{beaut} if
\begin{itemize}
\item $a$ divides $f(a)$ for all $a \in\{1,2, \ldots, n\}$, and
\item $f(a) \neq f(b)$ for all distinct elements $a$ and $b$
of $\{1,2, \ldots, n\}$ with $\operatorname{gcd}(a, b) \neq 1$.
\end{itemize}
Determine all positive integers $n$ such that, for every positive
integer $m$, there exists a beaut function
$f:\{1,2, \ldots, n\} \rightarrow\{m+1, m+2, \ldots, m+n\}$.
\end{problem}

\begin{solution}{Solution}
First we are going to show that for $n \geq 6$, a beaut function never
exists. \\

Suppose $n$ is even and $n \geq 6$. Pick $m = n(n-2) - 3$ so that the
image of $f$ is
$$
\{ n(n-2) - 2,\ n(n-2) - 1,\ n(n-2), \ldots, n(n-2) + (n-3) \}
$$
If we reduce this set modulo $n$ or $(n-2)$ it becomes
$$
\{ -2, \ -1, \ 0, \ 1, \ 2, \ldots, (n-3) \}
$$
So $n(n-2)$ is the only number that is divisible by $n$ or $(n-2)$.
Now if a beaut function $f$ existed we would require
$f(n) = f(n-2) = n(n-2)$ due to the divisiblity condition of beaut
functions. However $\gcd(n, n-2)$ is a multiple of 2 since both $n, (n-2)$
are even. But this contradicts the second condition of beaut functions.\\

Similarly suppose $n$ is odd and $n \geq 7$. Pick $m = (n-1)(n-3) - 4$ so
that the image of $f$ is
$$
\{ (n-1)(n-3) -3, \ (n-1)(n-3) -2, \ (n-1)(n-3) -1, \ (n-1)(n-3),
\ldots,
(n-1)(n-3) + (n-4) \}
$$
Reducing this set modulo $(n-1)$ or $(n-3)$ we have
$$
\{ -3, \ -2, \ -1, \ 0, \ 1, \ldots, (n-4) \}
$$
So $(n-1)(n-3)$ is the only number divisible by $(n-1)$ or $(n-3)$.
If a beaut function $f$ existed we would need $f(n-1) = f(n-3) = (n-1)(n-3)$
due to the divisibility condition of beaut functions. However
$\gcd(n-1, n-3)$ is even since both $(n-1), (n-3)$ are even. This
contradicts the second condition of beaut functions. \\

It remains to show that beaut functions do exist for $n=1,2,3,4,5$.
The $n=1$ case is trivial since 1 divides every integer. For the remaining
$n$, note that any set of $n$ consecutive numbers will have all the
residues mod $n$, so we will be able to find numbers in the image of
$f$ such that $a \in \{1,2,\ldots, n\}$ divides $f(a)$. Since $2,4$
is the only pair of non-coprime integers here we just need to ensure
$f(2) \neq f(4)$. This is possible since in any set of $\geq 4$ numbers
there should be one number $\equiv 0$ mod 4 and one number $\equiv 2$ mod 4.
$\blacksquare$
\end{solution}

\end{document}
